% Chapter 6: Discussion
\chapter{Discussion}
\label{ch:discussion}

This chapter critically reflects on the design decisions, trade-offs, and
limitations of OrderVista.  Each section evaluates an alternative that was
considered and justifies the approach taken, satisfying the assessment
criterion of demonstrating \emph{appropriate evaluation of alternative
approaches in problem solving}.

% ─────────────────────────────────────────────────────────────
\section{Route Optimisation: Nearest-Neighbour vs.\ Advanced Solvers}
\label{sec:disc-routing}

The nearest-neighbour (NN) heuristic was chosen as the sequencing algorithm
(Section~\ref{sec:route-impl}).  Table~\ref{tab:route-eval} showed a
consistent 25--31\% distance reduction over unoptimised ordering, confirming
that the heuristic adds genuine value.  However, NN is a greedy algorithm
that does not revisit earlier decisions; it can produce routes that are
10--25\% longer than the global optimum for medium-scale instances
\autocite{laporte2009fifty}.

Three stronger alternatives were evaluated during design:

\begin{enumerate}
  \item \textbf{2-opt local search}: iteratively reverses sub-sequences to
        remove crossing edges.  Typical improvement over NN is 5--10\%, with
        $O(n^{2})$ per iteration.  This would have been feasible for
        batch sizes up to $\sim$50 stops and is the most natural next step.
  \item \textbf{Ant Colony Optimisation (ACO)}: probabilistic
        metaheuristic that converges toward near-optimal tours over hundreds
        of iterations \autocite{dorigo1997ant}.  The computational cost
        (seconds to minutes) conflicts with the sub-second response
        requirement for mid-shift re-sequencing on a mobile device.
  \item \textbf{Google OR-Tools / OSRM}: server-side solvers that accept
        distance matrices and return optimised routes.  These require a
        persistent network connection, which contradicts the offline-first
        requirement, and introduce an additional external dependency.
\end{enumerate}

The decision to ship NN first was driven by two practical constraints:
(a)~mid-shift re-sequencing must execute on-device in under 200\,ms so
drivers are not blocked, and
(b)~the offline-first architecture precludes reliance on server-side solvers.
The modular design of the route pipeline (Section~\ref{sec:route-impl})
ensures that a 2-opt post-processing pass can be added without altering the
surrounding workflow. This is identified as the highest-priority future
improvement.

% ─────────────────────────────────────────────────────────────
\section{Backend-as-a-Service vs.\ Custom Server}
\label{sec:disc-baas}

Supabase was selected over a custom Node.js/Express backend and over Firebase
(Section~\ref{sec:tech-stack}).  In retrospect, this decision yielded three
significant benefits and one notable limitation:

\paragraph{Benefits.}
\begin{itemize}
  \item \textbf{Development velocity}: authentication, real-time WebSocket
        subscriptions, and file storage were available out of the box,
        eliminating weeks of boilerplate server code.
  \item \textbf{Row-Level Security}: PostgreSQL RLS policies enforce
        multi-tenant isolation at the database engine level, removing the risk
        of application-layer access-control bugs
        (Section~\ref{sec:security-test}).
  \item \textbf{SQL expressiveness}: complex queries (e.g.\ the
        fleet-wide assignment view joining tasks, drivers, and haversine
        distances) were written as standard SQL rather than the limited
        query syntax of Firestore.
\end{itemize}

\paragraph{Limitation.}
Supabase Edge Functions (Deno-based) were not used because the project
required Dart-only logic for consistency with the Flutter codebase.  All
business logic therefore resides in the client, which means that a
compromised client could, in theory, submit invalid state transitions
(e.g.\ marking a task as \texttt{delivered} without a PoD photo).
While RLS prevents cross-tenant access, it does not enforce workflow
constraints.  A production deployment would benefit from PostgreSQL
\texttt{BEFORE UPDATE} triggers or server-side functions to validate state
transitions. This is acknowledged as a limitation.

% ─────────────────────────────────────────────────────────────
\section{AI Model Selection: DeepSeek R1 vs.\ Alternatives}
\label{sec:disc-ai}

DeepSeek~R1 was selected for summarisation, conversational assistance, and
agentic action recommendation (Section~\ref{sec:ai-agentic}).  The decision
was evaluated against three alternatives:

\begin{table}[H]
  \centering
  \caption{AI model comparison for on-device logistics assistance.}
  \label{tab:ai-comparison}
  \small
  \begin{tabularx}{\textwidth}{l X X X X}
    \toprule
    \textbf{Criterion} & \textbf{DeepSeek R1} & \textbf{GPT-4o} &
    \textbf{Gemini 1.5} & \textbf{Local SLM} \\
    \midrule
    Cost per 1M tokens & \$0.55 & \$5.00 & \$3.50 & Free \\
    Structured JSON     & Reliable & Reliable & Moderate & Poor \\
    Latency (p50)       & $\sim$1.2\,s & $\sim$0.8\,s & $\sim$1.0\,s &
                          $\sim$3\,s \\
    Privacy             & Cloud & Cloud & Cloud & On-device \\
    Context window      & 64k & 128k & 1M & 4--8k \\
    \bottomrule
  \end{tabularx}
\end{table}

GPT-4o offers marginally lower latency but at roughly 9$\times$ the cost,
which would be prohibitive for a startup-stage fleet management product.
On-device small language models (SLMs) eliminate cloud dependency but
currently lack the structured-output reliability needed for agentic action
parsing; the 6/6 accuracy demonstrated in Table~\ref{tab:agentic-eval}
depends on consistent JSON output.  DeepSeek~R1 via OpenRouter provided the
best balance of cost, accuracy, and structured-output fidelity for this
project.

A notable risk is vendor lock-in: if OpenRouter or DeepSeek discontinue the
model, the prompt templates would need minor adaptation for a replacement LLM.
The system prompt and JSON schema are isolated in a single service class
(\texttt{DeepSeekService}), limiting the blast radius of such a change.

% ─────────────────────────────────────────────────────────────
\section{Offline-First Architecture: Trade-offs}
\label{sec:disc-offline}

The Drift-based local SQLite mirror (Section~\ref{sec:offline-sync}) ensures
that drivers can complete deliveries without connectivity.  The exponential
backoff sync queue successfully uploaded all queued items in unit testing
(Section~\ref{sec:unit-tests}).  However, two edge cases remain unresolved:

\begin{enumerate}
  \item \textbf{Conflict resolution}: if an admin reassigns a task while the
        driver is offline and simultaneously marking it as delivered, the
        current strategy is \emph{last-write-wins} based on the
        \texttt{updated\_at} timestamp.  A production system would benefit
        from operational-transform or CRDT-based merging.
  \item \textbf{Stale cache invalidation}: the local cache is refreshed on
        app launch and on WebSocket reconnection, but there is no
        time-to-live (TTL) eviction.  A driver who remains offline for
        several hours may act on outdated task data.  Adding a TTL-based
        warning banner is a low-effort improvement.
\end{enumerate}

Despite these limitations, the architecture proved sufficient for the
single-tester evaluation scenario and demonstrates the core offline-first
pattern that would underpin a production deployment.

% ─────────────────────────────────────────────────────────────
\section{Multi-Tenant Security: Strengths and Gaps}
\label{sec:disc-security}

The RLS-based isolation model passed all cross-tenant tests
(Section~\ref{sec:rls-test}), and the 17 policies across 11 tables
represent a defence-in-depth posture unusual for an academic prototype.
Two gaps are acknowledged:

\begin{itemize}
  \item \textbf{No penetration testing}: the security verification was
        limited to functional RLS queries executed through the Supabase
        client.  A professional penetration test (e.g.\ OWASP ZAP scan of
        the PostgREST API) was not performed.
  \item \textbf{Client-side secret storage}: API keys for Mapbox and
        OpenRouter are embedded in the Flutter build.  While obfuscated by
        the Dart AOT compiler, they could be extracted from the APK.  A
        production deployment should proxy these calls through a
        server-side gateway.
\end{itemize}

% ─────────────────────────────────────────────────────────────
\section{Comparison with Commercial Platforms}
\label{sec:disc-commercial}

Table~\ref{tab:commercial-compare} positions OrderVista against three
established commercial last-mile platforms.

\begin{table}[H]
  \centering
  \caption{Feature comparison with commercial last-mile platforms.}
  \label{tab:commercial-compare}
  \small
  \begin{tabularx}{\textwidth}{X c c c c}
    \toprule
    \textbf{Feature} & \textbf{OrderVista} & \textbf{Onfleet} &
    \textbf{Route4Me} & \textbf{Circuit} \\
    \midrule
    Route optimisation            & \checkmark & \checkmark & \checkmark & \checkmark \\
    Multi-tenant isolation        & \checkmark & \checkmark & ---        & ---        \\
    AI task summarisation         & \checkmark & ---        & ---        & ---        \\
    Agentic AI (action execution) & \checkmark & ---        & ---        & ---        \\
    Offline delivery completion   & \checkmark & Partial    & ---        & \checkmark \\
    Photographic PoD              & \checkmark & \checkmark & \checkmark & \checkmark \\
    Real-time admin--driver chat  & \checkmark & ---        & ---        & ---        \\
    Open-source / self-hostable   & \checkmark & ---        & ---        & ---        \\
    Dynamic priority scoring      & \checkmark & ---        & ---        & ---        \\
    Fleet-wide assignment         & \checkmark & \checkmark & \checkmark & ---        \\
    \bottomrule
  \end{tabularx}
\end{table}

OrderVista's differentiators are its \textbf{agentic AI layer}, its
\textbf{real-time admin--driver communication channel}, and its
\textbf{open-source, self-hostable architecture}.  Commercial platforms
offer more mature routing engines (often backed by proprietary solvers)
and broader hardware integration (barcode scanners, e-signatures).
OrderVista is not production-competitive on routing quality, but it
demonstrates capabilities, particularly agentic AI and dynamic
priority, that are absent from all three commercial comparators at the
time of writing.  Emerging industry evidence supports the promise of this
direction: one case study reports that an agentic AI framework for logistics
(autonomously reassigning drivers and rerouting vehicles) produced a 35\%
reduction in delivery time \autocite{maxinfocom2025agentic}, suggesting
that the pattern explored in OrderVista has significant real-world potential.

% ─────────────────────────────────────────────────────────────
\section{Limitations and Threats to Validity}
\label{sec:disc-limitations}

The following limitations should be considered when interpreting the
results of this project:

\begin{enumerate}
  \item \textbf{Single-tester evaluation}: all functional tests were
        performed by the developer.  Independent user testing with real
        dispatchers and drivers would provide stronger usability evidence.
  \item \textbf{Simulated workloads}: the route-optimisation evaluation
        used synthetic Hong Kong addresses rather than live delivery
        manifests.  Real-world routes include time-window constraints,
        vehicle capacity, and road restrictions not modelled here.
  \item \textbf{No load testing}: the system was not subjected to
        concurrent multi-user stress testing.  Supabase's connection
        pooling and PostgreSQL's MVCC provide theoretical scalability, but
        this was not empirically validated.
  \item \textbf{Android-only deployment}: the iOS build was not
        tested on a physical device due to the absence of an Apple
        Developer account.  Flutter's cross-platform guarantee mitigates
        but does not eliminate the risk of platform-specific issues.
  \item \textbf{AI non-determinism}: despite low temperature settings,
        LLM responses are inherently stochastic.  The 6/6 agentic accuracy
        reported in Table~\ref{tab:agentic-eval} may not generalise to
        all possible task configurations.
\end{enumerate}

% ─────────────────────────────────────────────────────────────
\section{Reflection on Development Process}
\label{sec:disc-process}

The project was developed iteratively across five phases
(Section~\ref{sec:methodology}).  Two reflections are notable:

\begin{itemize}
  \item \textbf{Incremental complexity}: beginning with core CRUD and
        adding AI features last proved effective.  The stable task pipeline
        provided a reliable test bed for the more experimental agentic
        layer.  Had agentic AI been attempted first, the lack of
        underlying data infrastructure would have made evaluation
        impossible.
  \item \textbf{BaaS productivity}: Supabase reduced the backend
        development effort to approximately 20\% of what a custom server
        would have required.  The time saved was reinvested into the
        AI integration and the fleet-optimisation algorithms, which are
        the project's primary contributions.  This confirms the
        literature finding that BaaS platforms accelerate prototyping
        at the cost of server-side flexibility
        \autocite{supabase2026docs}.
\end{itemize}
